\selectlanguage{english}
\frontmattertitle{Abstract}

\noindent
Coronary artery disease (CAD) assessment from coronary computed tomography angiography (CCTA) remains largely manual and time-consuming in clinical practice. At Hospital de la Santa Creu i Sant Pau, expert review of a single study can require between 30 minutes and two hours depending on case complexity. This bachelor's thesis presents an automated, end-to-end computational pipeline that transforms segmented coronary anatomy into quantitative stenosis metrics, CAD-RADS~2.0-oriented reporting support, and an interactive clinical visualization dashboard. The system is organized into four modular blocks: automated anatomy extraction, geometric stenosis quantification, segment labeling with CAD-RADS scoring, and a Streamlit-based decision-support interface. Designed as a transparent in-house framework rather than a proprietary black box, the pipeline aims to reduce reporting burden, improve reproducibility, and assist patient prioritization while preserving final clinical authority. Validation combines ASOCA and MACS-18 cohorts with synthetic phantoms for controlled geometric verification.

\selectlanguage{catalan}
\vspace{2em}
\frontmattertitle{Resum}

\noindent
L'avaluació de la malaltia coronària (MAC) a partir de la tomografia computada coronària (TCCC) continua sent, en la pràctica clínica, un procés manual i molt costós en temps. A l'Hospital de la Santa Creu i Sant Pau, la revisió experta d'un estudi pot requerir entre 30 minuts i dues hores segons la complexitat del cas. Aquest treball de fi de grau presenta un pipeline computacional automatitzat d'extrem a extrem que transforma l'anatomia coronària segmentada en mètriques quantitatives d'estenosi, suport orientat al CAD-RADS~2.0 i un tauler clínic interactiu de visualització. El sistema s'organitza en quatre blocs modulares: extracció automatitzada de l'anatomia, quantificació geomètrica de l'estenosi, etiquetatge segmentari amb puntuació CAD-RADS i una interfície Streamlit de suport a la decisió. Dissenyat com a marc transparent i intern, el pipeline pretén reduir la càrrega de informes, millorar la reproduïbilitat i assistir la priorització de pacients sense substituir l'autoritat clínica final. La validació combina les cohorts ASOCA i MACS-18 amb fantasmes sintètics per a la verificació geomètrica controlada.

\selectlanguage{spanish}
\vspace{2em}
\frontmattertitle{Resumen}

\noindent
La evaluación de la enfermedad coronaria (EC) a partir de la angiografía por tomografía computarizada coronaria (ATCC) sigue siendo, en la práctica clínica, un proceso manual y muy costoso en tiempo. En el Hospital de la Santa Creu i Sant Pau, la revisión experta de un estudio puede requerir entre 30 minutos y dos horas según la complejidad del caso. Este trabajo de fin de grado presenta un \emph{pipeline} computacional automatizado de extremo a extremo que transforma la anatomía coronaria segmentada en métricas cuantitativas de estenosis, soporte orientado al CAD-RADS~2.0 y un panel clínico interactivo de visualización. El sistema se organiza en cuatro bloques modulares: extracción automatizada de la anatomía, cuantificación geométrica de la estenosis, etiquetado segmentario con puntuación CAD-RADS e interfaz Streamlit de apoyo a la decisión. Diseñado como un marco transparente e interno, el \emph{pipeline} pretende reducir la carga de informes, mejorar la reproducibilidad y asistir la priorización de pacientes sin sustituir la autoridad clínica final. La validación combina las cohortes ASOCA y MACS-18 con fantomas sintéticos para la verificación geométrica controlada.

\selectlanguage{english}
